\documentclass[10pt]{article}
\providecommand{\FIGDIR}{fig}
\newcommand{\DIGESTDATE}{5 August 2026}
% ===================== PM Research Watch - shared preamble =====================
\usepackage[T1]{fontenc}
\usepackage[utf8]{inputenc}
\usepackage{amsmath,amssymb}
\usepackage{tgpagella}
\usepackage{tgheros}
\renewcommand{\sfdefault}{qhv}
\renewcommand{\ttdefault}{lmtt}
\usepackage{microtype}
\usepackage[english]{babel}

\usepackage[a4paper,top=17mm,bottom=17mm,left=17mm,right=17mm,
            headsep=6mm,footskip=10mm]{geometry}
\usepackage{graphicx}
\usepackage{xcolor}
\usepackage{colortbl}
\usepackage{tikz}
\usetikzlibrary{calc}
\usepackage[most]{tcolorbox}
\usepackage{titlesec}
\usepackage{fancyhdr}
\usepackage{booktabs}
\usepackage{array}
\usepackage{longtable}
\usepackage{enumitem}
\usepackage{multicol}
\usepackage{ragged2e}
\usepackage[hidelinks]{hyperref}

% ------------------------------------------------------------------ palette
\definecolor{Ink}    {HTML}{14262E}
\definecolor{Deep}   {HTML}{0F4C5C}
\definecolor{Teal}   {HTML}{1F7A8C}
\definecolor{Sky}    {HTML}{4FA3C4}
\definecolor{Sage}   {HTML}{6FA287}
\definecolor{Amber}  {HTML}{E3A008}
\definecolor{Clay}   {HTML}{B5643C}
\definecolor{Coral}  {HTML}{D1495B}
\definecolor{Violet} {HTML}{6A5B9E}
\definecolor{Slate}  {HTML}{4A6572}
\definecolor{Paper}  {HTML}{FBF9F5}
\definecolor{Mist}   {HTML}{EEEAE2}
\definecolor{Rule}   {HTML}{D9D3C9}

\pagecolor{Paper}
\color{Ink}

\hypersetup{colorlinks=true, linkcolor=Deep, urlcolor=Teal, citecolor=Teal,
            pdftitle={PM Research Watch}, pdfauthor={Automated daily digest}}

% ------------------------------------------------------------------ headings
\titleformat{\section}
  {\sffamily\bfseries\large\color{Deep}}{}{0pt}
  {}[{\vspace{1.2mm}\color{Amber}\titlerule[1.1pt]}]
\titlespacing*{\section}{0pt}{5.5mm}{2.8mm}

\titleformat{\subsection}
  {\sffamily\bfseries\normalsize\color{Teal}}{}{0pt}{}
\titlespacing*{\subsection}{0pt}{3.6mm}{1.4mm}

\setlist[itemize]{leftmargin=4.2mm, itemsep=1.1mm, topsep=1.3mm, parsep=0pt,
                  label=\textcolor{Teal}{\footnotesize$\blacktriangleright$}}

% ------------------------------------------------------------------ header/footer
\pagestyle{fancy}
\fancyhf{}
\renewcommand{\headrulewidth}{0pt}
\renewcommand{\footrulewidth}{0pt}
\fancyhead[L]{\sffamily\scriptsize\color{Slate}PM RESEARCH WATCH\;\textcolor{Rule}{\textbar}\;\DIGESTDATE}
\fancyhead[R]{\sffamily\scriptsize\color{Slate}Particulate matter monitoring \& health effects}
\fancyfoot[C]{\sffamily\scriptsize\color{Slate}\thepage}

% ------------------------------------------------------------------ components
% metric tile
\newcommand{\metric}[3]{%
  \begin{tcolorbox}[enhanced, colback=white, colframe=Rule, boxrule=0.5pt,
    arc=1.6mm, left=2.2mm, right=2.2mm, top=1.8mm, bottom=1.8mm,
    borderline west={1.7mm}{0pt}{#3}, nobeforeafter, valign=center,
    height=17mm, width=\linewidth]
    \RaggedRight
    {\sffamily\bfseries\fontsize{17}{17}\selectfont\color{#3}#1}\\[0.5mm]
    {\sffamily\fontsize{6.5}{7.4}\selectfont\color{Slate}#2}
  \end{tcolorbox}}

% pull-quote / key insight box
\newtcolorbox{keybox}[1][]{enhanced, colback=Mist, colframe=Mist,
  arc=1.6mm, left=3.2mm, right=3.2mm, top=2.4mm, bottom=2.4mm,
  borderline west={1.7mm}{0pt}{Amber}, #1}

% subtopic band
\newcommand{\band}[2]{%
  \vspace{3.2mm}\noindent
  \begin{tcolorbox}[enhanced, colback=#2, colframe=#2, arc=1.2mm,
    left=2.6mm, right=2.6mm, top=1.1mm, bottom=1.1mm, boxrule=0pt,
    width=\linewidth, nobeforeafter]
    {\sffamily\bfseries\small\color{white}#1}
  \end{tcolorbox}\vspace{1.6mm}}

% one paper entry: {short cite}{title}{journal}{doi}{body}
\newcommand{\paperentry}[5]{%
  \par\vspace{1.9mm}\noindent
  {\sffamily\bfseries\footnotesize\color{Ink}#1}\;%
  {\sffamily\fontsize{7}{8}\selectfont\color{Slate}#3}\\[0.4mm]
  {\footnotesize\itshape\color{Deep}#2}\\[0.6mm]
  {\footnotesize #4}\\[0.3mm]
  {\fontsize{6.6}{7.6}\selectfont\color{Slate}\texttt{doi:}\,\href{https://doi.org/#5}{\texttt{\detokenize{#5}}}}
  \par}

% caption style
\newcommand{\figcap}[1]{\par\vspace{1.0mm}
  {\fontsize{7.2}{8.6}\selectfont\color{Slate}\textsf{\textbf{Fig.}}\; #1\par}\vspace{1.5mm}}

\providecommand{\permil}{\textperthousand}

\setlength{\parindent}{0pt}
\setlength{\parskip}{0pt}
\renewcommand{\arraystretch}{1.18}


\begin{document}

% ==================================================================== MASTHEAD
\thispagestyle{empty}
\begin{tikzpicture}[remember picture, overlay]
  \fill[Deep] (current page.north west) rectangle
      ($(current page.north east)+(0,-67mm)$);
  \fill[Amber] ($(current page.north west)+(0,-67mm)$) rectangle
      ($(current page.north east)+(0,-68.4mm)$);
  \foreach \x/\y/\r in {22/12/0.9, 41/26/1.5, 63/9/0.7, 88/31/1.1, 112/17/1.9,
                        138/29/0.8, 159/13/1.3, 178/54/1.0, 196/21/1.6,
                        30/44/0.6, 74/57/1.2, 125/48/0.9, 168/6/0.8, 191/59/0.7,
                        52/18/1.0, 100/60/0.6, 148/20/1.1, 66/49/1.4, 205/45/0.8}
    \fill[Sky, opacity=0.30] ($(current page.north west)+(\x mm,-\y mm)$) circle (\r mm);
\end{tikzpicture}

\vspace*{4mm}
{\sffamily\fontsize{7.6}{9}\selectfont\color{Sky}%
 AUTOMATED DAILY LITERATURE DIGEST\;\textbullet\;PhD RESEARCH WATCH\par}
\vspace{2.2mm}
{\sffamily\bfseries\fontsize{27}{29}\selectfont\color{white}%
 Particulate Matter\par}
{\sffamily\fontsize{27}{29}\selectfont\color{Sky}%
 Monitoring \& Health Effects\par}
\vspace{5.0mm}
{\sffamily\fontsize{8.4}{10}\selectfont\color{white}%
 Issue 2026-08-05\;\;\textcolor{Amber}{\textbar}\;\;Entry window 5 August 2026\;\;%
 \textcolor{Amber}{\textbar}\;\;10 records in scope\par}
\vspace{18mm}

% -------------------------------------------------------------- metric strip
\noindent
\begin{minipage}[t]{0.190\linewidth}\metric{10}{RECORDS IN SCOPE\\(23 rejected, 7 held)}{Deep}\end{minipage}\hfill
\begin{minipage}[t]{0.190\linewidth}\metric{5}{SENSING RECORDS\\(50\% OF CORPUS)}{Teal}\end{minipage}\hfill
\begin{minipage}[t]{0.190\linewidth}\metric{4}{SUBTOPIC CLUSTERS\\REPRESENTED}{Sky}\end{minipage}\hfill
\begin{minipage}[t]{0.190\linewidth}\metric{6}{TIER-A\\RECORDS}{Amber}\end{minipage}\hfill
\begin{minipage}[t]{0.190\linewidth}\metric{1}{EFFECT ESTIMATE\\EXTRACTED}{Clay}\end{minipage}

\vspace{6mm}

% -------------------------------------------------------------- signal of day
\section*{\sffamily\bfseries\color{Deep}Signal of the day}
\vspace{-1.5mm}{\color{Amber}\rule{\linewidth}{1.1pt}}\vspace{2.5mm}

\begin{keybox}
{\sffamily\bfseries\small\color{Deep}Monitoring capacity is now measurably
\emph{anti}-correlated with need: in eastern Canada, situationally vulnerable populations
had 54\% higher odds of living in a PM$_{2.5}$ monitoring desert after adjusting for the
concentration itself.}\\[1.4mm]
\footnotesize
Siu et al.\ (\textit{Environ Res Health}) applied \textbf{locally adaptive kernel density
estimation} to all PM$_{2.5}$ monitors --- regulatory and low-cost --- across Ontario,
Quebec and Atlantic Canada, 2020--2024, setting each bandwidth from the local variability
of surface PM$_{2.5}$ and its sensitivity to distinguishing an exceedance. That is the
methodological move worth stealing: coverage is defined against the concentration field's
\emph{own} spatial structure rather than against a fixed radius, so ``monitored'' means
``able to resolve an exceedance here'', not ``within $x$\,km of a device''. In
mixed-effects logistic regression, residential instability carried \textbf{25\% higher
odds} and situationally vulnerable populations \textbf{54\% higher odds} of monitoring
deprivation, \emph{after} controlling for surface PM$_{2.5}$ and emission-facility
density --- so this is not simply monitors following pollution. With Canada's annual
standard tightening to 8\,$\mu$g\,m$^{-3}$ in 2030, the exceedance question becomes
harder exactly where the network is thinnest. Read alongside Jafari et al.\ and the
siting papers of 4 August, a coherent strand is forming: \textbf{siting is an
optimisation problem with an objective function}, and the objective is not coverage area
but resolvable exposure contrast in the populations that carry the burden.
\end{keybox}

\vspace{3mm}

% -------------------------------------------------------------- exec brief
\section*{\sffamily\bfseries\color{Deep}What changed today}
\vspace{-1.5mm}{\color{Amber}\rule{\linewidth}{1.1pt}}\vspace{2.5mm}

\footnotesize
% 3 lines of the last bullet spilled onto an otherwise blank p.2 on the first proof.
% Pulling the page down by 3 baselines absorbs them; PIPELINE 2026-08-04 records that
% 4-7 baselines starts printing over the folio, so this stays small and the bottom
% margin was re-proofed after the change.
\enlargethispage{3\baselineskip}
\begin{itemize}\setlength{\itemsep}{1.1mm}
\item \textbf{Entry window 5 August 2026, contiguous with the 4 August issue; no gap.}
The health axis was the weakest this watch has recorded: PubMed returned 15 health-axis
and 8 sensing-axis records for the 5 August entry date, of which \textbf{11 were rejected
as off-topic} --- diet scores, dementia checklists, a lead-exposure oncology paper, a
correction notice, therapeutic aerosols, poultry bioaerosols. The signal came from the
non-PubMed legs.

\item \textbf{Half the corpus is sensing, and for the second consecutive issue that is
network \emph{design}, not calibration chemistry.} Siu et al.\ quantify where monitoring
is absent; Jafari et al.\ solve fixed-and-mobile placement jointly as a constrained
optimisation over a Rotterdam bus network, maximising spatiotemporal coverage subject to
a budget \emph{and} to a constraint that every mobile sensor passes a reference station
often enough to stay calibrated. That last constraint is the interesting one: it makes
recalibration cadence a term in the siting objective rather than a separate operational
concern.

\item \textbf{One held DOI from 4 August was recovered and shipped.} Jafari et al.\
(IEEE OJ-ITS) had no retrievable abstract in Crossref, Europe PMC or Consensus on 4
August; a Semantic Scholar Graph query returned it on 5 August. \textbf{Semantic Scholar
is now a fourth abstract source in the retry chain} --- it is the only index that resolved
any of the five held records. The other four, all Elsevier, remain unretrievable after a
second attempt and stay held rather than rejected.

\item \textbf{G\"abel \& Hertig answer the question the calibration literature keeps
deferring: how often.} A yearlong Augsburg co-location across $-$10 to 36\,$^{\circ}$C
and 19--96\,\% RH finds that a single pre-deployment calibration decays, and that
\textbf{monthly in-season recalibration} is what keeps an SPS30 inside the EU/EPA
indicative-measurement data-quality objective. Carried as an out-of-window backfill
(\textit{Atmos.\ Meas.\ Tech.}, created 19 February 2026) --- \textit{AMT} is structurally
invisible to PubMed, which is exactly the gap the Consensus leg exists to close. It must
be excluded from any weekly trend line.

\item \textbf{Sea breeze as a monitoring-representativeness problem, not a meteorology
result.} Yang et al.\ identify 428 diurnal sea-breeze events in a 10-year coastal record
and show that excluding them changes the diurnal amplitude \emph{and} the seasonal cycle
of O$_3$ and CH$_4$ under southwesterly flow. A coastal background site is therefore not
a background site on those days. Aerosol number showed a pronounced diurnal cycle
attributed to new particle formation from accumulated ship SO$_2$ --- a mechanism that
generates ultrafine number with negligible mass, invisible to a mass-based network.

\item \textbf{Only one effect estimate carries an interval, and it is from a preprint.}
Zhang et al.\ (medRxiv, not peer reviewed) estimate 7.1\% (95\% UI 5.4--8.8\%) of
cardiorespiratory ED visits averted by combined EGU, mobile-source and port policies
across NYC, LA and Atlanta. Fourth consecutive issue without a peer-reviewed PM ratio
estimate with a confidence interval --- see the forest-plot caption so a rollup does not
read this as pipeline loss.
\end{itemize}

\clearpage

% ==================================================================== ANALYTICS
\section*{\sffamily\bfseries\color{Deep}Corpus analytics}
\vspace{-1.5mm}{\color{Amber}\rule{\linewidth}{1.1pt}}\vspace{3mm}

\noindent\includegraphics[width=\linewidth]{\FIGDIR/f1_subtopics.png}
\figcap{Subtopic assignment of the 10 in-scope records, single-label by dominant
contribution. \textit{Sensing, forecasting \& instrumentation} takes 5 of 10 (50\%) --- the
highest share this watch has recorded --- ahead of exposure assessment and burden/policy at
two each. \textbf{Six of the ten standing subtopics are empty}, including every
mechanistic, respiratory and reproductive cluster: an unusually narrow window, driven by
an exceptionally poor health-axis yield rather than by screening.}

\vspace{2mm}
\noindent\includegraphics[width=\linewidth]{\FIGDIR/f2_design_endpoint.png}
\figcap{Left --- study architecture. Modelling, inventory and geospatial work leads with
5/10, ahead of measurement campaigns at 3; there is \textbf{no experimental toxicology
and no cohort epidemiology in this window at all}, and the single acute-observational
entry is a preprint. Right --- health endpoint by organ system. Seven of ten records carry
no human health endpoint, consistent with a corpus that is half instrumentation. Note ---
the design map fell through to \textit{Other / mixed} for six records on the first pass
of this run, the fourth recurrence of that class; the cause was canonical group names
being written directly onto records, and the map is now closed under its own output (see
run log).}

\vspace{2mm}
\noindent\includegraphics[width=\linewidth]{\FIGDIR/f3_metric_geography.png}
\figcap{Left --- particle metric. Four records use PM$_{2.5}$ as the exposure; the rest are
unspeciated PM, PM source factors from positive matrix factorisation, precursor-nitrogen
emissions, or paired PM$_{2.5}$/PM$_{10}$/number. Right --- geographic provenance. North
America and Europe carry three each and China contributes only one --- an inversion of the
last two issues, where China supplied nearly half the corpus.}

\vspace{2mm}
\noindent\includegraphics[width=\linewidth]{\FIGDIR/f6_lifecourse.png}
\figcap{Life-course windows addressed; counts are non-exclusive. Only four of ten records
carry a human life-course window at all. The childhood entry is the South Texas school
deployment; the older-adult entry is a burden projection rather than a measured cohort,
so the apparent coverage is thinner than the bar heights suggest.}

\vspace{2mm}
\noindent\includegraphics[width=\linewidth]{\FIGDIR/f4_heatmap.png}
\figcap{Subtopic $\times$ study-architecture cross-tabulation. The standing structural
feature persists: the \textit{Sensing} row contains no health-endpoint linkage. With 10
records across four subtopics the matrix is very sparse --- read it as a coverage map, not
a distribution.}

\vspace{2mm}
\noindent\includegraphics[width=\linewidth]{\FIGDIR/f5_forest.png}
\figcap{\textbf{The single estimate plotted is a policy-attributable fraction from a
preprint, not a peer-reviewed PM concentration--response estimate.} Zhang et al.\
(medRxiv) report 7.1\% (95\% UI 5.4--8.8\%) of cardiorespiratory ED visits averted by
combined EGU, mobile-source and port policies across three US cities, 2005--2019, with
median PM$_{2.5}$ reductions of 27--62\%. No peer-reviewed record in this entry window
reported a PM$_{2.5}$ or PM$_{10}$ ratio estimate with a confidence interval --- the fourth
consecutive issue for which that holds. The plot is included so the absence is recorded
explicitly rather than inferred from a missing figure.}

\clearpage

% ==================================================================== DIGEST
\section*{\sffamily\bfseries\color{Deep}Paper-level digest}
\vspace{-1.5mm}{\color{Amber}\rule{\linewidth}{1.1pt}}\vspace{1mm}

{\fontsize{7.4}{9}\selectfont\color{Slate}Ordered by cluster. Each entry gives the
methodological core and the finding that matters, not an abstract paraphrase.\par}

\band{Sensing, forecasting \& instrumentation\hfill 5 records}{Amber}

\paperentry{Siu et al.\ 2026}
{Disparities in fine particulate matter (PM$_{2.5}$) monitoring capacity in eastern
Canada}
{Environ Res Health}
{Locally adaptive kernel density estimation of total PM$_{2.5}$ monitor density ---
regulatory plus low-cost --- across Ontario, Quebec and Atlantic Canada, 2020--2024, with
each local bandwidth set from the surrounding variability of surface PM$_{2.5}$ and its
sensitivity to resolving an exceedance. Mixed-effects logistic regression gives
\textbf{25\% higher odds of monitoring deprivation for residential instability and 54\%
higher odds for situationally vulnerable populations}, after adjusting for surface
PM$_{2.5}$ and emission-facility density; suburban Montreal and Toronto, the Windsor and
Niagara border regions, southeast Kitchener and remote northwest Ontario are named.
Limitation: monitor density is a proxy for information yield, and low-cost and regulatory
devices are pooled without weighting by data quality, so a dense but uncalibrated cluster
counts as capacity. The adaptive-bandwidth construction is directly reusable for
designing where to add nodes rather than merely auditing where they are.}
{10.1088/2752-5309/ae87c0}

\paperentry{Jafari et al.\ 2026}
{Joint optimization of bus-based mobile sensing and reference sensor deployment for urban
air quality monitoring}
{IEEE Open J Intell Transp Syst}
{Mathematical-programming formulation selecting simultaneously \textbf{which buses to
instrument with low-cost sensors and where to place reference stations}, maximising
spatiotemporal network coverage under a budget constraint while requiring that every
mobile sensor pass a reference site frequently enough to be recalibrated. Demonstrated on
a Rotterdam case study covering residential and industrial areas. The contribution is
structural rather than empirical: recalibration cadence enters the siting objective as a
constraint, so coverage and data quality are traded off in one problem instead of
sequentially. Limitation: coverage is geometric --- optimised against route and population
geometry, with no validation against a measured concentration field, and no treatment of
the sampling bias that fixed bus routes impose. This is the record held on 4 August and
recovered on 5 August via Semantic Scholar.}
{10.1109/ojits.2026.3706855}

\paperentry{Yang et al.\ 2026}
{Diurnal sea breeze worsens coastal air quality and complicates monitoring of remote
North Atlantic air}
{Atmos Chem Phys}
{Ten years of observations from Penlee Point Atmospheric Observatory on the northeast
Atlantic coast, from which \textbf{428 diurnal sea-breeze events} were identified, most
frequent in spring and summer under low wind and strong irradiance. Daytime surface
O$_3$, NO$_x$, CH$_4$, CO$_2$, SO$_2$, total aerosol number, PM$_{2.5}$ and PM$_{10}$ were
all elevated during events, with mean daytime O$_x$ near \textbf{45\,ppb} and increased
regulatory exceedance for O$_3$. Accumulated ship-emitted SO$_2$ appears to drive new
particle formation and a pronounced diurnal cycle in aerosol number. The monitoring
payload: \textbf{excluding sea-breeze events reduces the diurnal amplitude and alters the
seasonal cycle of O$_3$ and CH$_4$} under southwesterly flow --- a coastal ``background''
site is conditionally not background. Limitation: one site, so the event classification
and its consequences are not shown to transfer. Number-without-mass is precisely the
regime a mass-based low-cost network cannot see.}
{10.5194/acp-26-10909-2026}

\paperentry{G\"abel \& Hertig 2026}
{Recalibration of low-cost O$_3$ and PM$_{2.5}$ sensors: linking practices to recent air
sensor test protocols}
{Atmos Meas Tech}
{Yearlong urban-background co-location at Augsburg of Alphasense OX-B431 and Sensirion
SPS30 against reference instruments, spanning $-$10 to 36\,$^{\circ}$C, 19--96\,\% RH and
PM$_{2.5}$ to 153\,$\mu$g\,m$^{-3}$, with MLR, ridge, random forest and XGBoost
calibration functions evaluated against the current US EPA and CEN air-sensor test
protocols. A five-month baseline training gives PM$_{2.5}$ $R^2$ 0.93--0.97 in training
and \textbf{0.84--0.93 in the following test period}, RMSE 2.04--3.02\,$\mu$g\,m$^{-3}$
--- the decay is the finding. Testing pairwise recalibration cycles, \textbf{monthly
in-season recalibration} was required to keep uncertainty inside the data-quality
objective for indicative measurements. Limitation: one site, one year, two sensor models,
so the cadence is an existence result rather than a general rule. \textbf{Out-of-window
backfill} (created 19 Feb 2026); \textit{AMT} is not indexed by PubMed. Exclude from
weekly trend lines.}
{10.5194/amt-19-1293-2026}

\paperentry{Wang et al.\ 2026}
{Traffic-induced environmental health assessment at intersections using roadside LiDAR}
{Environ Monit Assess}
{Roadside LiDAR extracts high-resolution vehicle trajectories at one urban and one
suburban intersection, from which traffic emissions, noise and ground vibration are
estimated and combined into a fuzzy comprehensive assessment score, with weights set by
quadratic programming under ranking constraints, variable-weight grey entropy, and game
theory to reconcile the two. The urban intersection scored \textbf{69.45 by day, falling
to 57.12 at night ($\sim$18\%)}, while the suburban site was stable at 72.80/74.54;
population-dimension scores were consistently below environmental-dimension scores.
Limitation: PM is \emph{inferred} from trajectories through an emission model and never
measured, so this is a modelled exposure surrogate validated against nothing --- the
elaborate weighting machinery sits on top of an unvalidated emissions estimate. Included
because trajectory-derived emission proxies are an increasingly common input to near-road
sensor fusion, and their weak point should be explicit.}
{10.1007/s10661-026-15750-y}

\clearpage

\band{Exposure assessment \& modelling\hfill 2 records}{Sky}

\paperentry{Pinakana et al.\ 2026}
{Influence of transboundary wildfire smoke on fine particulate matter concentrations at
South Texas border schools}
{NPJ Clean Air}
{Low-cost sensors deployed at schools in Roma and Rio Grande City, Starr County, on the
US--Mexico border for \textbf{146 days, November 2023 to April 2024}, corrected by
collocation-based calibration against a research-grade instrument and read alongside
meteorology, back trajectories and satellite smoke and fire products. Daily mean
PM$_{2.5}$ was significantly higher on smoke-impacted days ($p<0.001$), with
\textbf{site-level increases of 8.61\% to 54.9\%} --- a sixfold spread across sites within
one small county, which is itself the finding: transboundary smoke impact is not
spatially uniform at the scale a single regulatory monitor would represent. Trajectories
attribute the source to Mexican wildfire emissions. Limitation: one season, no indoor
measurement despite the school framing, and attribution rests on trajectory plus
satellite coincidence rather than on composition. A clean template for low-cost sensor
deployment in an under-monitored community.}
{10.1038/s44407-026-00096-z}

\paperentry{Phalen et al.\ 2026}
{Inhaled aerosol dosimetry update}
{Inhal Toxicol}
{Synthesis of the Fourth International Aerosol Dosimetry Conference (Irvine, 2024)
covering computational deposition models, in vitro and in vivo exposure systems, and
deposition variability. Reported advances: integration of model types to span all major
respiratory-tract regions; improved in vitro particle delivery and dose measurement; the
upper airways recognised as a significant entry route. Two findings matter for ambient PM
work --- \textbf{dose variability emerged as the critical cross-cutting issue}, and
\textbf{deposited ultrafine particles may acquire surface coatings after deposition and
translocate into tissue, including fetal tissue}. Limitation: a conference report, not a
systematic review --- coverage reflects who presented, there are no effect sizes and no
pooling, and the fetal-translocation claim is stated without the primary evidence. Useful
mainly as a statement of where dosimetric uncertainty sits, which is the term that
converts a measured ambient concentration into an internal dose.}
{10.1080/08958378.2026.2691193}

\band{Burden, policy \& mitigation\hfill 2 records}{Teal}

\paperentry{Zhang et al.\ 2026 (preprint)}
{Health impacts of national and local air pollution control policies targeting electric
generating units, mobile sources, and port activities in three US cities}
{medRxiv (preprint, not peer reviewed)}
{Counterfactual policy evaluation for New York City, Los Angeles and Atlanta, 2005--2019.
Daily PM$_{2.5}$, criteria gases, 1-in-3-day PM$_{2.5}$ components and PMF-resolved PM
sources at six monitoring sites were used to construct the concentrations that would have
occurred without the selected EGU, mobile-source and port policies; city-site-specific
multi-pollutant Poisson models then linked observed concentrations to respiratory
emergency-department visits and hospitalised cardiovascular ED visits, and the fitted
models were applied to both concentration series. \textbf{Median PM$_{2.5}$ reductions of
27--62\%} across cities, and \textbf{7.1\% (95\% UI 5.4--8.8\%) of cardiorespiratory
visits averted}. Limitation: not peer reviewed; the counterfactual is a modelled construct
whose uncertainty is not propagated into the health interval, and multi-pollutant Poisson
models over highly correlated criteria pollutants cannot cleanly assign the benefit to PM.
The source-resolved PMF layer is the part worth following.}
{10.64898/2026.08.04.26359640}

\paperentry{Goyal et al.\ 2026}
{Quantifying environmental co-benefits of nitrogen-based crop restructuring and its
implications on India's interstate trade network}
{Nat Commun}
{Integrated nutrient, water, greenhouse-gas, economic and interstate-trade accounting for
Indian agriculture, comparing nitrogen-focused against water-focused crop-restructuring
pathways. The nitrogen-focused pathway cuts \textbf{nitrogen surplus by 13.4\%, water use
by 18.6\% and greenhouse gases by 8.7\%} while holding baseline calorie production, and
delivers about \textbf{4.6-fold larger reciprocal co-benefits} than the water-focused
alternative, worth roughly \textbf{USD 1.19\,bn} in avoided nitrogen-pollution
socio-environmental cost. Limitation: PM is entirely implicit --- reactive nitrogen is a
dominant precursor of secondary inorganic aerosol over the Indo-Gangetic Plain, but no
ambient PM$_{2.5}$ response is modelled and no health endpoint is estimated, so the
PM-relevant co-benefit is asserted through the nitrogen accounting rather than computed.
Carried because agricultural NH$_3$ is the least-constrained term in South Asian PM
source apportionment.}
{10.1038/s41467-026-75905-w}

\band{Cardiovascular \& metabolic\hfill 1 record}{Clay}

\paperentry{Shi et al.\ 2026}
{Ischemic heart disease in G20 countries: forecasted burden and regional disparities
based on GBD 2023 estimates}
{Catheter Cardiovasc Interv}
{GBD 2023 descriptive burden analysis of ischemic heart disease incidence, mortality and
DALYs across the G20, 1990--2023, with EAPC and joinpoint regression, population
attributable fractions for modifiable risks, ARIMA and Bayesian age-period-cohort
projection to 2050, and XGBoost with SHAP for DALY determinants. \textbf{8.52\,M cases
and 5.81\,M deaths in 2023}; age-standardised incidence and mortality fell (EAPC $-$1.61
and $-$1.49) but incidence rose in women aged 15--49, with Russia highest on all three
age-standardised rates and Japan projected to peak again around 2032. High systolic blood
pressure dominated the risk attribution, with \textbf{particulate matter and lead named as
the prominent environmental contributors}. Limitation: GBD population attributable
fractions are modelled from the same exposure--response literature this watch tracks, so
this is not independent evidence of a PM effect --- it is that literature propagated
through a burden model, and the machine-learning layer adds interpretability, not
identification. Tier C: read as context, not as evidence.}
{10.1002/ccd.70783}

\clearpage

% ==================================================================== TABLE
\section*{\sffamily\bfseries\color{Deep}Corpus register}
\vspace{-1.5mm}{\color{Amber}\rule{\linewidth}{1.1pt}}\vspace{2.5mm}

{\fontsize{6.7}{8.4}\selectfont
\setlength{\tabcolsep}{2.4pt}
\begin{longtable}{@{}>{\RaggedRight}p{25mm}>{\RaggedRight}p{41mm}>{\RaggedRight}p{27mm}%
>{\RaggedRight}p{23mm}>{\RaggedRight}p{18mm}>{\RaggedRight}p{15mm}@{}}
\toprule
\textsf{\textbf{First author}} & \textsf{\textbf{Focus}} & \textsf{\textbf{Journal}} &
\textsf{\textbf{Design}} & \textsf{\textbf{Metric}} & \textsf{\textbf{Region}}\\
\midrule
\endfirsthead
\toprule
\textsf{\textbf{First author}} & \textsf{\textbf{Focus}} & \textsf{\textbf{Journal}} &
\textsf{\textbf{Design}} & \textsf{\textbf{Metric}} & \textsf{\textbf{Region}}\\
\midrule
\endhead
\multicolumn{6}{@{}l@{}}{\cellcolor{Mist}\textsf{\textbf{\textcolor{Deep}{Neuro / mental health}}}}\\[0.6mm]
\href{https://doi.org/10.1001/jamanetworkopen.2026.25015}{Dales} & Wildfire-sourced PM$_{2.5}$ and ED visits for migraine and other primary headache & \textit{JAMA Netw Open} & Case-crossover & PM$_{2.5}$ (wildfire) & Canada\\
\href{https://doi.org/10.1093/ageing/afag215}{Glover} & Umbrella meta-review of urban environmental exposures and cognitive health & \textit{Age Ageing} & Umbrella review / meta-analysis & PM$_{2.5}$ & Global\\
\href{https://doi.org/10.1016/j.nsa.2026.107019}{Goossens} & In utero air pollution exposure and pre-/postnatal brain development & \textit{Neurosci Appl} & Narrative review & PM$_{2.5}$ / PM$_{10}$ / PM$_{\mathrm{coarse}}$ & Global\\
\href{https://doi.org/10.1016/j.isci.2026.116617}{Ma} & PM$_{2.5}$-induced neurotoxicity and depression: exposure risks and mechanisms & \textit{iScience} & Narrative review & PM$_{2.5}$ + components & Global\\
\href{https://doi.org/10.3760/cma.j.cn112150-20250804-00752}{Pan} & Maternal PM exposure and offspring neurodevelopment: effects and mechanisms & \textit{Zhonghua Yu Fang Yi Xue Za Zhi} & Narrative review & PM (bulk) & China\\
\addlinespace[0.9mm]
\multicolumn{6}{@{}l@{}}{\cellcolor{Mist}\textsf{\textbf{\textcolor{Teal}{Cardiovascular \& metabolic}}}}\\[0.6mm]
\href{https://doi.org/10.3760/cma.j.cn112338-20251120-00843}{Huang} & PM$_{2.5}$ and dyslipidaemia in rural Chinese children; dietary effect modification & \textit{Zhonghua Liu Xing Bing Xue Za Zhi} & Cross-sectional (multistage) & PM$_{2.5}$ & China\\
\href{https://doi.org/10.32394/pe/218966}{Kazmierska} & Air pollution and cardiovascular disease: review of current data & \textit{Przegl Epidemiol} & Narrative review & PM$_{2.5}$ / PM$_{10}$ & Global\\
\href{https://doi.org/10.1007/s00125-026-06814-2}{Xi} & Air pollution, built environment and temperature vs. prediabetes phenotypes & \textit{Diabetologia} & Prospective cohort & PM$_{2.5}$ & Germany\\
\href{https://doi.org/10.1161/JAHA.125.046917}{Yang} & Short-term PM$_{2.5}$ and hospitalisation/ED use in atrial fibrillation & \textit{J Am Heart Assoc} & Case-crossover & PM$_{2.5}$ (wildfire-influenced) & USA\\
\addlinespace[0.9mm]
\multicolumn{6}{@{}l@{}}{\cellcolor{Mist}\textsf{\textbf{\textcolor{Sky}{Reproductive \& developmental}}}}\\[0.6mm]
\href{https://doi.org/10.1016/j.jes.2025.10.043}{Chen} & Prenatal air pollution, fetal growth trajectory phenotypes and child neurodevelopment & \textit{J Environ Sci (China)} & Prospective cohort & PM$_{2.5}$ / PM$_{10}$ & China\\
\href{https://doi.org/10.1016/j.ecoenv.2026.120566}{Huang} & Joint prenatal bisphenol + air-pollutant mixture exposure and preschool IQ & \textit{Ecotoxicol Environ Saf} & Prospective cohort (mixtures) & PM$_{2.5}$ / PM$_{10}$ & China\\
\href{https://doi.org/10.1007/s40572-026-00552-8}{Peterson} & Ambient air pollution, PAHs and pubertal development: critical review & \textit{Curr Environ Health Rep} & Critical review & PM$_{2.5}$ components & Global\\
\href{https://doi.org/10.1016/j.envres.2026.125313}{Xie} & Long-term PM$_{2.5}$ and timing of pubertal development in boys and girls & \textit{Environ Res} & Prospective cohort & PM$_{2.5}$ & China\\
\addlinespace[0.9mm]
\multicolumn{6}{@{}l@{}}{\cellcolor{Mist}\textsf{\textbf{\textcolor{Sage}{Respiratory \& allergic}}}}\\[0.6mm]
\href{https://doi.org/10.1007/s00484-026-03289-z}{Bayar Muluk} & Climate change and otolaryngology: airway disease impacts and emission reduction & \textit{Int J Biometeorol} & Narrative review & PM$_{2.5}$ & Global\\
\href{https://doi.org/10.4103/mgr.MEDGASRES-D-25-00164}{Liao} & Short-term air pollution, respiratory symptoms and lung function in asthma & \textit{Med Gas Res} & Panel study & PM$_{2.5}$ / PM$_{10}$ & China\\
\href{https://doi.org/10.3390/epidemiologia7040097}{Sanchez de Prada} & Ambient PM and asthma-coded admissions in children 0-3 y, Spain & \textit{Epidemiologia} & Ecological panel & PM$_{2.5}$ / PM$_{10}$ & Spain\\
\href{https://doi.org/10.1016/j.jhazmat.2026.143076}{Yang} & Allergenic vs. toxic risk of airborne seed hairs compared with ambient PM$_{2.5}$ & \textit{J Hazard Mater} & Field measurement & PM$_{2.5}$ + bioaerosol & China\\
\addlinespace[0.9mm]
\multicolumn{6}{@{}l@{}}{\cellcolor{Mist}\textsf{\textbf{\textcolor{Violet}{Mechanistic toxicology}}}}\\[0.6mm]
\href{https://doi.org/10.1016/j.jes.2025.12.074}{Ding} & PM$_{2.5}$ promotes vascular calcification via miR-27b-5p/HES1 signalling & \textit{J Environ Sci (China)} & In vitro + organoid + murine & PM$_{2.5}$ & China\\
\href{https://doi.org/10.1007/s00360-026-01698-5}{Maraschi} & Osmoregulatory disruption and metal bioaccumulation from settleable atmospheric PM & \textit{J Comp Physiol B} & Controlled exposure (ecotox) & Settleable APM & Brazil\\
\href{https://doi.org/10.1016/j.jes.2025.11.018}{Wu} & Macrophage NLRP3 inflammasome in real-ambient PM-induced hepatic glucose dysregulation & \textit{J Environ Sci (China)} & Real-ambient murine + co-culture & Real-ambient PM & China\\
\addlinespace[0.9mm]
\multicolumn{6}{@{}l@{}}{\cellcolor{Mist}\textsf{\textbf{\textcolor{Clay}{Occupational \& indoor}}}}\\[0.6mm]
\href{https://doi.org/10.3389/fmedt.2026.1796828}{Chen} & Nested suction instrument to cut intraoperative surgical-smoke PM$_{2.5}$ & \textit{Front Med Technol} & Bench / device experiment & PM$_{2.5}$ (surgical smoke) & China\\
\href{https://doi.org/10.1093/annweh/wxag055}{Cole} & Determinants of respirable crystalline silica exposure in Australian tunnelling & \textit{Ann Work Expo Health} & Observational exposure study & Respirable dust / RCS & Australia\\
\href{https://doi.org/10.1016/j.jhazmat.2026.143081}{Ding} & Stage-resolved hair Zn-isotope dynamics as a firefighting combustion-exposure tracer & \textit{J Hazard Mater} & Repeated-measures biomarker & Combustion-derived PM & China\\
\href{https://doi.org/10.3390/jox16040119}{Ramsden} & Xenobiotic hazards in aircraft cabin air & \textit{J Xenobiot} & Review + risk economics & Ultrafine particles & Global\\
\addlinespace[0.9mm]
\multicolumn{6}{@{}l@{}}{\cellcolor{Mist}\textsf{\textbf{\textcolor{Amber}{Sensing, forecasting \& instrumentation}}}}\\[0.6mm]
\href{https://doi.org/10.1016/j.jenvman.2026.130515}{Shi} & Adaptive spatiotemporal graph transformer for multi-site multi-step PM$_{2.5}$ forecasting & \textit{J Environ Manage} & Deep-learning model & PM$_{2.5}$ & China / India\\
\href{https://doi.org/10.3390/bios16070380}{Tsuchiya} & Seasonal and indoor-outdoor characteristics of fluorescent aerosol particles in offices & \textit{Biosensors} & Real-time sensor campaign & Bioaerosol / coarse & Japan\\
\addlinespace[0.9mm]
\multicolumn{6}{@{}l@{}}{\cellcolor{Mist}\textsf{\textbf{\textcolor{Amber}{Exposure assessment \& modelling}}}}\\[0.6mm]
\href{https://doi.org/10.1029/2025GH001621}{Solorzano-Araque} & Multiscale exposure modelling: LOTOS-EUROS + mobility surveys + particle toxicity & \textit{GeoHealth} & CTM + mobility integration & PM$_{2.5}$ / PM$_{10}$ & Colombia\\
\addlinespace[0.9mm]
\multicolumn{6}{@{}l@{}}{\cellcolor{Mist}\textsf{\textbf{\textcolor{Coral}{Burden, policy \& mitigation}}}}\\[0.6mm]
\href{https://doi.org/10.1002/hsr2.72886}{Karamzad} & Ischaemic heart disease burden in MENA 1990-2021; PM attributable fraction 28.2\% & \textit{Health Sci Rep} & GBD secondary analysis & Ambient PM & MENA\\
\href{https://doi.org/10.1371/journal.pone.0347989}{Rees} & HEPA filtration to reduce respiratory infection in care homes & \textit{PLoS One} & RCT process evaluation & Filtered PM & UK\\
\href{https://doi.org/10.1016/j.jes.2026.01.009}{Shan} & Policy-driven emission reductions and post-pandemic rebound, high-resolution inventories & \textit{J Environ Sci (China)} & Emission inventory & PM emissions & China\\
\addlinespace[0.9mm]
\multicolumn{6}{@{}l@{}}{\cellcolor{Mist}\textsf{\textbf{\textcolor{Slate}{Other clinical endpoints}}}}\\[0.6mm]
\href{https://doi.org/10.3390/dj14070440}{Lupita} & PM$_{2.5}$, preventive-care capacity and caries/edentulism burden across the EU & \textit{Dent J} & Ecological panel & PM$_{2.5}$ & EU-27\\
\href{https://doi.org/10.1016/j.jdin.2026.06.008}{Park} & Long-term PM exposure and risk of androgenetic alopecia & \textit{JAAD Int} & Nationwide cohort & PM$_{2.5}$ / PM$_{10}$ & South Korea\\
\href{https://doi.org/10.3171/2026.3.PEDS2615}{Sun} & Environmental exposure and area deprivation vs. paediatric hydrocephalus outcomes & \textit{J Neurosurg Pediatr} & Registry cohort + geospatial & PM$_{2.5}$ (satellite-derived) & USA\\
\addlinespace[0.9mm]
\bottomrule
\end{longtable}}


\vspace{4mm}
\begin{keybox}
{\sffamily\bfseries\footnotesize\color{Deep}Provenance and method}\\[1.2mm]
{\fontsize{7.2}{8.8}\selectfont
Entry window \textbf{5 August 2026}, contiguous with the 4 August issue; no gap. Nine of
the ten shipped records carry a 5 August entry or Crossref \texttt{created} date; the
tenth (G\"abel \& Hertig) is an explicitly flagged out-of-window backfill.
\textbf{Sources.} PubMed E-utilities and the PubMed connector on two separate axes
(health, instrumentation) windowed on \texttt{[EDAT]} --- 15 health and 8 sensing hits
after merging the two, of which 11 were rejected as off-topic. Europe PMC
\texttt{CREATION\_DATE} --- 2 hits, 1 carried (a medRxiv preprint). Crossref scoped by
journal ISSN across eight aerosol and atmospheric titles that PubMed does not index ---
7 hits, 1 carried, 3 rejected as out of scope, 3 held for want of an abstract. Consensus
on three sensing axes (calibration and drift, satellite--surface fusion, carbonaceous
source apportionment) --- every hit resolved to a DOI created before this window;
\textbf{1 backfill carried, 9 rejected}. OpenAlex returned HTTP 429 and remains
paid-blocked for date windowing; arXiv returned nothing in scope after the recency screen.
\textbf{Screening.} Deduplication keys PMID $\rightarrow$ lowercased DOI $\rightarrow$
title signature, checked against \texttt{seen.json} before summarising; Consensus titles
were resolved to DOIs through Crossref \emph{before} screening, which is what prevented
three re-summaries of records already shipped. \textbf{Seven records are held, not
rejected}: their DOIs verify in Crossref but no abstract is retrievable from Crossref,
Europe PMC, Semantic Scholar or Consensus. \textbf{Semantic Scholar was added to the
retry chain on this run} and recovered one of the five records held on 4 August; the
remaining four, all Elsevier, stay held. Held records are deliberately absent from
\texttt{seen.json} so a later issue can carry them. Nothing was summarised from a title
alone, and no DOI was inferred from a journal prefix.
\textbf{Labels.} Subtopic, design, particle-metric and geography assignments are
single-label judgements from title, abstract, keywords and MeSH terms, not a controlled
vocabulary. The single effect estimate is transcribed verbatim from a preprint; see the
forest-plot caption.
All DOI links resolve to the publisher of record and were checked by the automated DOI
gate before this issue compiled. Abstracts remain the copyright of their publishers.\par}
\end{keybox}

\end{document}
